\chapter{Han \emph{Yuefu} and the Nineteen Old Poems}

\section{The Music Bureau}

\begin{definition}[\emph{Yuefu}]
\zht{樂府} \emph{yuèfǔ} literally means ``Music Bureau'', a government office (expanded
under Emperor Wu of Han, r.\ 141--87 BCE) that collected songs from the regions and
composed music for court ritual and entertainment. By extension \emph{yuefu} means the
\emph{songs} associated with it, and later any poem written to an old \emph{yuefu} title,
even when the music was long lost.
\end{definition}

\begin{important}[Why \emph{yuefu} matter]
\begin{itemize}
  \item They mark the move from the four-character line to the \textbf{five-character
        line} (\zht{五言}), which becomes the main line of classical \emph{shi} for the
        next thousand years. Many \emph{yuefu} also have irregular lines.
  \item They are \textbf{narrative and dramatic}: dialogue, characters, stories of ordinary
        people (soldiers, orphans, abandoned wives, a girl picking mulberry leaves).
  \item Later poets (Li Bai, Du Fu, Bai Juyi) wrote ``new'' \emph{yuefu} using the genre
        to comment on society.
\end{itemize}
\end{important}

\section{Reading Han \emph{yuefu}}

\subsection{``South of the River'' \zht{江南}}

\begin{center}
\begin{tabular}{p{0.3\textwidth}p{0.3\textwidth}p{0.3\textwidth}}
\toprule
Original & Pinyin & My translation \\
\midrule
\zht{江南可採蓮，} & jiāngnán kě cǎi lián & South of the River, good for picking lotus: \\
\zht{蓮葉何田田。} & lián yè hé tiántián & how the lotus leaves crowd the water! \\
\zht{魚戲蓮葉間。} & yú xì lián yè jiān & Fish play among the lotus leaves; \\
\zht{魚戲蓮葉東，} & yú xì lián yè dōng & fish play east of the lotus leaves, \\
\zht{魚戲蓮葉西，} & yú xì lián yè xī & fish play west of the lotus leaves, \\
\zht{魚戲蓮葉南，} & yú xì lián yè nán & fish play south of the lotus leaves, \\
\zht{魚戲蓮葉北。} & yú xì lián yè běi & fish play north of the lotus leaves. \\
\bottomrule
\end{tabular}
\end{center}

\textbf{Analysis.} Almost nothing ``happens'': the poem is pure song structure, a call and
four answering lines (probably sung by different voices) circling the compass. But
\zht{蓮} \emph{lián} ``lotus'' is a homophone of \zht{憐} \emph{lián} ``to love'', and fish
playing among lotus is an old image of lovers. The repetition imitates the darting fish
and the to-and-fro of a work song for girls on boats.

\subsection{``Heaven Above!'' \zht{上邪}}

A love oath, with irregular line lengths:

\begin{center}
\begin{tabular}{p{0.38\textwidth}p{0.52\textwidth}}
\toprule
\zht{上邪！} & Heaven above! \\
\zht{我欲與君相知，} & I want to be your love \\
\zht{長命無絕衰。} & for ever, never failing, never fading. \\
\zht{山無陵，} & When the mountains have no peaks, \\
\zht{江水為竭，} & when the rivers run dry, \\
\zht{冬雷震震，} & when thunder rolls in winter, \\
\zht{夏雨雪，} & when snow falls in summer, \\
\zht{天地合，} & when heaven and earth come together --- \\
\zht{乃敢與君絕！} & only then will I dare to part from you! \\
\bottomrule
\end{tabular}
\end{center}

\textbf{Analysis.} A catalogue of \emph{impossibilities} (\emph{adynata}), each one more
cosmic, leading to a climactic final line. Lines of 2, 3, 4 and 5 characters: the
irregularity is a sign of real song. The speaker is (by convention) a woman; the directness
contrasts sharply with the restraint of later literati poetry.

\subsection{Narrative \emph{yuefu}}

\begin{itemize}
  \item ``Mulberry by the Path'' (\zht{陌上桑}): the beautiful Luofu, picking mulberry
        leaves, rebuffs a passing governor by boasting about her (perhaps imaginary)
        husband. Comic and clever.
  \item ``Fighting South of the Wall'' (\zht{戰城南}): the dead lie unburied; the speaker
        asks the crows to cry for them before eating them. A grim anti-war song.
  \item ``Southeast the Peacock Flies'' (\zht{孔雀東南飛}): the longest early Chinese
        narrative poem (over 350 lines), a tragedy of a young couple forced apart by his
        mother, ending in double suicide.
  \item ``The Ballad of Mulan'' (\zht{木蘭詩}): Northern Dynasties, not Han, but the most
        famous \emph{yuefu}: a daughter takes her father's place in the army.
\end{itemize}

\section{The Nineteen Old Poems \zht{古詩十九首}}

A group of nineteen anonymous five-character poems from the later Eastern Han (probably
2nd century CE), first grouped together in the anthology \emph{Wen xuan} \zht{文選}
(c.\ 530). They are the first great body of literati five-character verse.

\begin{important}[Themes of the Nineteen Old Poems]
\begin{itemize}
  \item \textbf{Separation}: the wife at home, the traveller on the road.
  \item \textbf{Brevity of life}: ``Man's life between heaven and earth / is like a
        traveller passing by''; therefore seize the day (wine, fine clothes).
  \item \textbf{Frustrated ambition}: the unrecognised scholar seeking office in the
        capital.
  \item \textbf{Style}: plain diction, natural images, no ornament, emotional directness.
        Later critics called them ``the crown of five-character poetry''.
\end{itemize}
\end{important}

\subsection{Poem 1: ``On and on, going on and on'' \zht{行行重行行}}

\begin{center}
\begin{tabular}{p{0.3\textwidth}p{0.6\textwidth}}
\toprule
\zht{行行重行行，} & On and on, going on and on, \\
\zht{與君生別離。} & alive, and parted from you. \\
\zht{相去萬餘里，} & Ten thousand miles and more between us, \\
\zht{各在天一涯。} & each at one edge of the sky. \\
\zht{道路阻且長，} & The road is hard and long; \\
\zht{會面安可知？} & who knows when we will meet again? \\
\zht{胡馬依北風，} & The northern horse leans into the north wind, \\
\zht{越鳥巢南枝。} & the southern bird nests on the southern branch. \\
\zht{相去日已遠，} & Each day we are farther apart; \\
\zht{衣帶日已緩。} & each day my sash hangs looser. \\
\zht{浮雲蔽白日，} & Drifting clouds cover the bright sun; \\
\zht{遊子不顧返。} & the wanderer does not think of coming back. \\
\zht{思君令人老，} & Longing for you makes me old; \\
\zht{歲月忽已晚。} & the years have suddenly grown late. \\
\zht{棄捐勿復道，} & Put it aside, speak of it no more: \\
\zht{努力加餐飯。} & only try your best to eat well. \\
\bottomrule
\end{tabular}
\end{center}

\subsubsection{Analysis}
\begin{itemize}
  \item \textbf{Voice.} A woman left behind (some readers: a friend or a loyal minister).
  \item \textbf{Opening.} \zht{行行重行行}: four of five characters are \zht{行} ``go'',
        enacting the endless journey.
  \item \textbf{Parallel couplet.} ``Northern horse / north wind'' against ``southern bird
        / southern branch'': even animals cling to home, so why not you? (\emph{Hu} and
        \emph{Yue} are the far north and far south.)
  \item \textbf{Physical detail.} The loosening sash shows her wasting away, a body image
        that stands for grief without naming it.
  \item \textbf{Drifting clouds cover the sun}: the absent man is distracted (by another
        woman?), or, allegorically, slanderers obscure the ruler.
  \item \textbf{Ending.} The final turn is to self-restraint and care: ``eat well'' may be
        addressed to him or to herself. The poem refuses a dramatic climax; that understated
        close is characteristic.
\end{itemize}

\subsection{Poem 10: ``Far, far away, the Herd-Boy Star'' \zht{迢迢牽牛星}}

\begin{center}
\begin{tabular}{p{0.3\textwidth}p{0.6\textwidth}}
\toprule
\zht{迢迢牽牛星，} & Far, far away, the Herd-Boy Star; \\
\zht{皎皎河漢女。} & bright, bright, the Weaving Girl of the River of Heaven. \\
\zht{纖纖擢素手，} & Slender, slender, she lifts her pale hands; \\
\zht{札札弄機杼。} & \emph{click-clack} she works the shuttle of her loom. \\
\zht{終日不成章，} & All day long she cannot finish a pattern, \\
\zht{泣涕零如雨。} & her tears fall like rain. \\
\zht{河漢清且淺，} & The River of Heaven is clear and shallow; \\
\zht{相去復幾許？} & how far apart are they, after all? \\
\zht{盈盈一水間，} & Brimming, a single stretch of water between them: \\
\zht{脈脈不得語。} & gazing, gazing, they cannot speak. \\
\bottomrule
\end{tabular}
\end{center}

\subsubsection{Analysis}
\begin{itemize}
  \item \textbf{Myth.} The Herd Boy (Altair) and the Weaving Girl (Vega) are lovers
        separated by the Milky Way (``River of Heaven''), allowed to meet once a year on
        the seventh night of the seventh month (\zht{七夕}).
  \item \textbf{Reduplicatives.} Six of ten lines open with or contain doubled characters
        (\zht{迢迢}, \zht{皎皎}, \zht{纖纖}, \zht{札札}, \zht{盈盈}, \zht{脈脈}),
        inherited from the \emph{Shijing} and here used for rhythm and longing.
  \item \textbf{Irony.} The river is ``clear and shallow'': the barrier is tiny, which makes
        the separation more painful, not less.
\end{itemize}

\section{From Han to the Jian'an poets}

At the end of the Han, the Cao family (Cao Cao \zht{曹操}, his sons Cao Pi and Cao Zhi
\zht{曹植}) and their circle (the ``Jian'an'' period, 196--220) made five-character poetry
the prestige form of personal expression. Cao Cao's ``Short Song'' (\zht{短歌行}) opens:

\begin{center}
\zht{對酒當歌，人生幾何！} --- ``Facing wine, we should sing: how long is a man's life?''
\end{center}

\begin{aside}
The ``Jian'an spirit'' (\zht{建安風骨}, ``Jian'an wind and bone'') became a critical
byword for vigorous, unornamented poetry with real feeling, and Tang poets invoked it
against the ornate court style of the intervening centuries.
\end{aside}

\section*{Summary}
\begin{itemize}
  \item \emph{Yuefu}: Music Bureau songs; folk origin, narrative, often irregular lines;
        later a genre for social commentary.
  \item ``South of the River'' (song structure, lotus pun), ``Heaven Above!''
        (\emph{adynata} oath).
  \item Nineteen Old Poems: anonymous, Eastern Han, five-character; separation, transience,
        plain style.
  \item Five-character line becomes standard; Jian'an poets (the Cao family) make it the
        literati form.
\end{itemize}
